% Part: first-order-logic
% Chapter: beyond
% Section: introduction

\documentclass[../../../include/open-logic-section]{subfiles}

\begin{document}

\olfileid{fol}{byd}{int}

\olsection{ټوليزه کتنه}

د لومړۍ درجې منطق د پام وړ يوازينے منطقي نظام نۀ دے: د دې منطق ډېر
غځېدلي او بدل ډولونه شته۔ يو منطق په عام ډول د يوې ژبې صوري مشخصات،
اکثره، خو تل نۀ، د استنتاج يو نظام، او اکثره، خو تل نۀ، يوه ټاکل شوې
معنٰی پوهنه لري۔ خو د دې ويي تخنيکي کارونه يوه څرګنده پوښتنه راپورته
کوي: هغه منطقونه چې د لومړۍ درجې منطق نۀ دي، په وجداني يا فلسفي معنا
د ``منطق'' له ويي سره څه اړيکه لري؟ لاندې ټول بيان شوي نظامونه د يوه
يا بل ډول استدلال د مدل کولو دپاره جوړ شوي؛ آيا ويلے شو چې څه شي دوى
منطقي کوي؟

اسانه ځوابونه نۀ تر سترګو کېږي۔ د ``منطق'' ويے په بېلابېلو لارو او
بېلابېلو چاپېريالونو کښې کارول کېږي، او دا مفهوم، لکه د ``رښتيا''
مفهوم، له ډېرو فلسفي دريځونو څېړل شوے دے۔ د بېلګې په توګه، څوک ښايي
د منطقي استدلال موخه دا وګڼي چې معلومه کړي کومې ويناوې په لازمي ډول
رښتونې، له تجربې وړاندې رښتونې، د غيرمنطقي اصطلاحاتو له تفسير څخه
خپلواکې رښتونې، د خپلې بڼې له مخې رښتونې، يا د ژبني تړون له مخې
رښتونې دي؛ او له دغو مفکورو څخه هره يوه ډېر وضاحت غواړي۔ که پام يوازې
په رياضياتو کښې کارېدونکي منطق ته هم محدود کړو، د هغه د ساحې په هکله لږ
توافق شته۔ د بېلګې په توګه، رسل او وايټ‌هېډ په
\textit{پرېنسيپيا مېتېماتيکا} کښې هڅه وکړه چې رياضيات د منطق پر بنسټ،
د فرېګې له پيل کړې {\em منطقپالې} دود سره سم، جوړ کړي۔ د دوى منطقي
نظام د لوړو ډولونو د منطق يوه بڼه وه چې لاندې بيان شوي نظام ته ورته
وه۔ په پاى کښې اړ شول داسې بديهي اصول ورزيات کړي چې د ډېرو معيارونو
له مخې سوچه منطقي نۀ ښکاري (په ځانګړي ډول د لامتناهي والي بديهي اصل
او د راکمېدنې بديهي اصل)؛ خو بيا هم څوک دا دريځ لرلے شي چې د لوړې
درجې استدلال ځينې بڼې بايد منطقي ومنل شي۔ د دې پر خلاف، کواين، چې
هستپوهنه ئې ``قضيې'' د بحث د روا څيزونو په توګه نۀ مني، استدلال کوي
چې د دويمې او لوړې درجې منطق په اصل کښې د سټ تيوري ده چې د پسه پوست
ئې اغوستے؛ په نورو ټکو، هغه نظامونه چې پر پريديکاتونو کميت ټاکنه لري
سوچه منطقي نۀ دي۔

اوس دپاره غوره ده چې دا ډول فلسفي مسئلې بل وخت ته پرېږدو او لاندې
نظامونه يوازې د بېلابېلو ډولونو استدلال، که منطقي وي که بل څه، صوري
ايډيالې بڼې وګڼو۔

\end{document}
